\section{ITERAZIONE 2}

In questa iterazione il sistema viene esteso con il primo dei due moduli
algoritmici del progetto: lo \textbf{Smistatore Intelligente Multi-Reparto}, che
introduce un instradamento adattivo delle comande verso i reparti di
preparazione. L'iterazione comprende l'evoluzione dei diagrammi UML (con il nuovo
sottosistema di smistamento e i nuovi componenti di gestione), la descrizione e
l'analisi dell'algoritmo, la suite di test dedicata, le API esposte e una prima
rilevazione delle metriche statiche del codice.

% ============================================================================
\subsection{Casi d'uso implementati}

% Rinumerazione: nell'Iterazione 0 UC5a e' il caso base "Genera Stampe e Smistare".
% Il modulo algoritmico che lo realizza in forma avanzata e' un caso d'uso distinto
% (UC7), per evitare che lo stesso ID cambi significato tra iterazioni.
\begin{table}[H]
\centering
\caption{Casi d'uso sviluppati nell'Iterazione 2}
\label{tab:uc-iter2}
\begin{tabular}{lll}
\toprule
\textbf{ID} & \textbf{Titolo} & \textbf{Relazione} \\
\midrule
UC7 & Smistatore Intelligente Multi-Reparto & realizza UC5a in forma avanzata \\
\bottomrule
\end{tabular}
\end{table}

Il caso d'uso UC5a (\textit{Genera Stampe e Smistare}), descritto nell'Iterazione~0
e implementato nella sua forma base, viene qui realizzato in forma avanzata dal
modulo algoritmico UC7. Per non sovraccaricare lo stesso identificatore con due
significati, il modulo riceve un ID dedicato.

% ============================================================================
\subsection{Diagrammi UML}

I diagrammi seguenti aggiornano quelli dell'Iterazione~1 per riflettere
l'introduzione del sottosistema di smistamento e dei nuovi componenti di
gestione. Per ciascun diagramma vengono evidenziate le differenze rispetto alla
versione precedente.

\subsubsection{UML Component Diagram}
Rispetto all'Iterazione~1, l'architettura è stata ampliata e si articola ora in cinque sottosistemi principali: \textit{Database}, \textit{Application Logic}, \textit{Client}, \textit{Stampanti} e il nuovo sottosistema \textit{Smistamento}.

Il sottosistema \textit{Application Logic} è organizzato in sei componenti. Accanto a \texttt{GestioneOrdini}, \texttt{GestioneMenu}, \texttt{GestioneScorte} e \texttt{GestioneStampe} sono formalizzate due nuove viste logiche: \texttt{GestioneAllergeni} e \texttt{GestioneConfigurazione}. La separazione è concettuale, finalizzata a distinguere la sola lettura sugli allergeni dall'amministrazione del menù; a livello implementativo entrambe confluiscono nel modulo \texttt{routes/menu.js}, condividendo lo stesso accesso al database. Il diagramma rappresenta pertanto tali componenti come viste sulla medesima sorgente dati, senza introdurre interfacce di persistenza fisicamente distinte.

Nel sottosistema \textit{Client} le interfacce sono state aggiornate per riflettere questa separazione di responsabilità. L'interfaccia di scrittura \texttt{MenuMgmt} viene dismessa e sostituita dalla nuova interfaccia \texttt{ConfigurazioneMgmt}, riservata al componente \textit{AdminFrontEnd}, che assorbe la logica amministrativa relativa a inserimento, modifica ed eliminazione delle voci, impostazione delle scorte iniziali, configurazione delle stampanti e associazione degli allergeni. La \textit{CassaFrontEnd} mantiene l'interfaccia di sola lettura \texttt{MenuLettura} per la visualizzazione della griglia prodotti e ottiene accesso a \texttt{ScorteMgmt}, prima non disponibile. L'interfaccia \texttt{AllergeniMgmt}, esposta dalla vista logica \texttt{GestioneAllergeni}, è condivisa da entrambi i client in modalità di sola lettura.

Infine, il flusso di stampa ha subito una revisione architetturale. \texttt{GestioneStampe} non comunica più direttamente con il sottosistema \textit{Stampanti}: è stato introdotto il sottosistema \textit{Smistamento}, contenente il componente \texttt{SistemaSmistamento}, che funge da dispatcher instradando le richieste verso i componenti \texttt{Stampante} tramite l'interfaccia \texttt{StampeMgmt}. La comunicazione tra \texttt{GestioneStampe} e \texttt{SistemaSmistamento} avviene attraverso l'interfaccia \texttt{SmistamentoMgmt}. Questa soluzione rende l'instradamento verso i reparti più scalabile e disaccoppiato rispetto all'approccio precedente.

\begin{figure}[H]
    \centering
    \includegraphics[width=0.85\textwidth]{IT2_UML_Component_Diagram.png}
    \caption{Diagramma dei Componenti UML  Iterazione 2}
    \label{fig:uml_component}
\end{figure}

\subsubsection{UML Class Diagram Interface}

\texttt{OrdiniMgmt} espone le operazioni per la gestione operativa e il ciclo di
vita di una comanda: inserimento, rimozione e modifica quantitativa delle pietanze,
inserimento di note libere per la cucina, gestione delle specifiche di asporto,
applicazione di sconti, calcolo del resto e procedure di conferma o annullamento
dell'ordine.

% COERENZA: l'interfaccia di lettura si chiama MenuLettura.
\texttt{MenuLettura} funge da modulo di sola consultazione ad uso della cassa,
esponendo metodi per l'ottenimento del menù e il filtraggio delle voci per
categoria o per settore. L'intera logica amministrativa e di scrittura è invece
delegata alla nuova interfaccia \texttt{ConfigurazioneMgmt}, riservata al
backoffice, che espone inserimento, modifica ed eliminazione fisica delle pietanze,
impostazione delle scorte iniziali, configurazione e recupero delle stampanti di
reparto e associazione degli allergeni ai singoli piatti.

La nuova interfaccia \texttt{AllergeniMgmt} rende fruibili le informazioni sulle
intolleranze alimentari, con metodi per consultare e stampare il report degli
allergeni sia per una specifica pietanza sia a livello globale per l'intero menù,
su entrambi i front-end.

L'interfaccia \texttt{ScorteMgmt} è stata potenziata per garantire la consistenza
dell'inventario durante il servizio: oltre a verificare la disponibilità dei piatti,
espone i metodi per decrementare le scorte residue e per gestire i contatori
aggregati associati ai piatti a disponibilità limitata.

Infine, l'invio alle stampanti è stato disaccoppiato tramite l'interfaccia
\texttt{SmistamentoMgmt}, legata al rispettivo sottosistema: essa si occupa
dell'elaborazione ad alto livello dei flussi di stampa (generazione e
instradamento dei biglietti verso i reparti corretti, produzione della copia
cliente, notifica dei piatti pronti per l'asporto). A valle, l'interfaccia
\texttt{StampeMgmt} gestisce la comunicazione a basso livello con i dispositivi
fisici, esponendo i metodi per l'invio del file di stampa e per il controllo dello
stato hardware delle stampanti.

\begin{figure}[H]
    \centering
    \includegraphics[width=0.85\textwidth]{IT2_UML_Class_Diagram_Interface.png}
    \caption{Diagramma UML delle Classi per interfacce  Iterazione 2}
    \label{fig:uml_classes_interface}
\end{figure}

\subsubsection{UML Class Diagram per tipo di dato}
Il Class Diagram per tipo di dato rappresenta le entità del dominio e i relativi
attributi, in corrispondenza diretta con lo schema relazionale definito in
\texttt{database/init.sql}: nomi e attributi coincidono con quelli effettivi delle
tabelle, così che diagramma e database descrivano la stessa struttura. Le entità
sono raggruppate per ambiti logici.

\paragraph{Menù.} L'entità \texttt{Voce} costituisce il cuore del catalogo, con
\texttt{codice}, \texttt{nome}, \texttt{prezzo}, \texttt{categoria},
\texttt{settore\_visualizzazione}, \texttt{settore\_stampa}, \texttt{colore\_tasto},
\texttt{visibile}, \texttt{asportabile}, \texttt{modalita\_stampa},
\texttt{copia\_scontrino\_cliente} e i parametri usati dai moduli algoritmici:
\texttt{tempo\_preparazione}, \texttt{tempo\_riapprovvigionamento} e
\texttt{priorita\_voce}. Le pietanze aggregate (menù composti) sono modellate dalla
relazione \texttt{VoceComposizione} (\texttt{voce\_aggregata\_id},
\texttt{voce\_componente\_id}).

\paragraph{Disponibilità.} La classe \texttt{Scorta} (\texttt{voce\_id},
\texttt{quantita}, \texttt{soglia\_giallo}, \texttt{soglia\_rosso},
\texttt{attiva}) mantiene la disponibilità per singola voce, con le soglie di
allerta come propri attributi; \texttt{ScortaStorico} (\texttt{voce\_id},
\texttt{quantita}, \texttt{data\_rifornimento}, \texttt{note}) registra i
rifornimenti. Gli ingredienti condivisi sono gestiti da
\texttt{ContatoreAggregato} (\texttt{nome}, \texttt{quantita}) e dalla relazione
\texttt{VoceContatore} (\texttt{voce\_id}, \texttt{contatore\_id}).

\paragraph{Comande.} L'entità \texttt{Ordine} (\texttt{tavolo}, \texttt{stato},
\texttt{asporto}, \texttt{sconto}, \texttt{tipo\_sconto}, \texttt{totale},
\texttt{importo\_pagato}, \texttt{timestamp}, \texttt{utente\_id}) traccia stato e
dati fiscali; \texttt{OrdineRiga} (\texttt{ordine\_id}, \texttt{voce\_id},
\texttt{quantita}) mappa le singole quantità; \texttt{OrdineRigaNota}
(\texttt{riga\_id}, \texttt{numero\_porzione}, \texttt{testo},
\texttt{costo\_aggiuntivo}) gestisce le note per porzione e
\texttt{NotaPreimpostata} (\texttt{voce\_id}, \texttt{testo},
\texttt{costo\_aggiuntivo}) le varianti rapide.

\paragraph{Allergeni.} L'entità \texttt{Allergene} (\texttt{nome}, \texttt{descr})
è associata alle voci tramite la relazione \texttt{VoceAllergene}
(\texttt{voce\_id}, \texttt{allergene\_id}).

\paragraph{Stampa.} L'entità \texttt{Stampante} (\texttt{reparto}, \texttt{nome},
\texttt{indirizzo\_ip}, \texttt{porta}, \texttt{modello}, \texttt{primaria},
\texttt{attiva}, \texttt{stato}) descrive i dispositivi di reparto, mentre
\texttt{StampaEseguita} (\texttt{ordine\_id}, \texttt{stampante\_id},
\texttt{reparto}, \texttt{payload}, \texttt{esito}, \texttt{durata\_ms},
\texttt{timestamp}) costituisce l'audit log delle stampe inviate.

\paragraph{Configurazione.} L'entità \texttt{Configurazione} (\texttt{chiave},
\texttt{valore}) conserva i parametri runtime usati dai moduli algoritmici (es.\
soglie e finestre temporali).

\begin{figure}[H]
    \centering
    \includegraphics[width=0.85\textwidth]{IT2_UML_Class_Diagram_tipo_di_dato.png}
    \caption{Diagramma UML di Classe per tipo di dato  Iterazione 2}
    \label{fig:uml_classes_dato}
\end{figure}

\newpage
\subsubsection{UML Deployment Diagram}

Il nodo \textit{FesteDB} ospita il database relazionale e funge da sorgente dati
centrale. Il nodo centrale \textit{Web Server} ospita il core applicativo,
organizzato nei moduli di gestione (\texttt{GestioneOrdini}, \texttt{GestioneMenu},
\texttt{GestioneScorte}, \texttt{GestioneAllergeni}, \texttt{GestioneConfigurazione}
e \texttt{GestioneStampe}); allergeni e configurazione condividono il modulo
\texttt{routes/menu.js} e quindi lo stesso accesso al database. All'interno del nodo
si nota una dipendenza logica tra \texttt{GestioneMenu} e \texttt{GestioneScorte}.

% COERENZA: la cassa accede a MenuLettura (lettura), non a MenuMgmt.
Il nodo \textit{Client} ospita le applicazioni di frontend, \texttt{WebAppCassa} e
\texttt{WebAppBackOffice}. La comunicazione tra Client e Web Server è mediata dalle
interfacce aggiornate: la cassa accede a \texttt{OrdiniMgmt}, \texttt{MenuLettura},
\texttt{ScorteMgmt} e \texttt{AllergeniMgmt}; il backoffice si affida anch'esso ad
\texttt{AllergeniMgmt} e utilizza in modo esclusivo \texttt{ConfigurazioneMgmt} per
l'amministrazione centralizzata.

Infine, l'infrastruttura di stampa si divide su due nodi logici distinti. Il Web
Server comunica tramite l'interfaccia \texttt{SmistamentoMgmt} con il nuovo nodo
\textit{Smistamento}, che ospita il componente \texttt{SistemaSmistamento}
delegato alla logica di routing. Questo nodo comunica a sua volta con i dispositivi
fisici allocati nel nodo \textit{Stampante} tramite l'interfaccia a basso livello
\texttt{StampeMgmt}.

\begin{figure}[H]
    \centering
    \includegraphics[width=0.85\textwidth]{IT2_UML_Deployment_Diagram.png}
    \caption{Diagramma UML di Deployment  Iterazione 2}
    \label{fig:uml_deployment}
\end{figure}

% ============================================================================
\newpage
\subsection{Algoritmo: Smistatore Intelligente Multi-Reparto}

\subsubsection*{Introduzione}

L'algoritmo \textbf{Smistatore Intelligente Multi-Reparto} bilancia il carico
operativo tra i reparti di un sistema di gestione ordini per la ristorazione.

Nei sistemi POS tradizionali, ogni voce d'ordine viene instradata verso un reparto
fisso (\texttt{settore\_stampa}), indipendentemente dal carico corrente. Questo
genera due problemi:

\begin{itemize}
    \item \textbf{Saturazione asimmetrica:} tutti gli ordini di un reparto
    confluiscono su di esso anche quando è sovraccarico.
    \item \textbf{Mancanza di adattività:} il sistema ignora il tempo di
    preparazione residuo, la quantità richiesta, lo stato delle stampanti e la
    disponibilità delle scorte.
\end{itemize}

L'algoritmo introduce un modello di scoring dinamico che valuta, per ogni task, il
reparto ottimale compatibile in base al carico reale istantaneo
(Figura~\ref{fig:flusso_smistatore}).

\paragraph{Attori coinvolti}

\begin{itemize}
    \item \textbf{Sistema POS:} emette gli ordini con i relativi parametri
    (timestamp, priorità, quantità).
    \item \textbf{Reparti operativi:} cucina, bar, griglia destinatari delle
    comande.
    \item \textbf{Stampanti / Dashboard:} destinatari fisici; il sistema gestisce
    fallback se offline.
    \item \textbf{Cassa:} riceve notifiche di blocco per scorta insufficiente.
\end{itemize}

\paragraph{Trigger e postcondizioni}

\begin{description}
    \item[Trigger:] Ricezione di un ordine con una o più righe (pietanze) e
    relativi metadati.
    \item[Postcondizione:] Per ogni riga valida: assegnazione al reparto con carico
    minore compatibile, generazione della comanda separata per reparto, invio alla
    stampante o dashboard corretta.
\end{description}

\subsubsection*{Parametri di Input}

Ogni riga d'ordine viene convertita in un \textit{task di preparazione}. La
Tabella~\ref{tab:parametri} riporta i parametri considerati per routing e ranking.

\begin{table}[H]
\centering
\caption{Parametri di input dell'algoritmo}
\label{tab:parametri}
\begin{tabular}{>{\ttfamily}lp{9cm}}
\toprule
\normalfont\textbf{Parametro} & \textbf{Descrizione} \\
\midrule
settore\_stampa        & Reparto primario configurato nel menu; punto di partenza per i reparti compatibili. \\
numero\_task\_in\_coda & Task in attesa per ogni reparto. Componente principale del carico. \\
tempo\_preparazione    & Tempo medio stimato di preparazione (secondi). \\
quantità               & Porzioni richieste. Incide sul carico e sul raggruppamento batch. \\
timestamp              & Marca temporale di arrivo; usata per urgenza e ordinamento. \\
stato\_stampante       & Online/offline. Attiva il fallback se offline. \\
disponibilità\_scorte  & Flag o quantità residua. Se insufficiente, il task viene bloccato. \\
priorità\_ordine       & Categoria: tavolo, asporto, urgenza. Peso nel ranking $P$. \\
\bottomrule
\end{tabular}
\end{table}

\subsubsection*{Passi dell'Algoritmo}

L'algoritmo si articola in undici passi sequenziali, con diramazioni per la gestione
dei casi limite. L'obiettivo è minimizzare il carico massimo tra i reparti e
massimizzare la reattività.

\paragraph{Passo 1  Ricezione e lettura righe ordine}
All'arrivo di un ordine, il sistema estrae da ogni riga il \texttt{settore\_stampa},
il timestamp, la priorità e la quantità. Ogni riga diventa un task di preparazione
indipendente.

\paragraph{Passo 2  Conversione in task di preparazione}
Ogni riga viene normalizzata in un oggetto task con struttura uniforme,
indipendentemente dal tipo di pietanza o reparto.

\paragraph{Passo 3  Verifica disponibilità scorte}
Prima del routing, il sistema verifica che la scorta sia sufficiente per la quantità
richiesta. Questo \textit{guard} evita di occupare capacità di reparto inutilmente.
Se la scorta è insufficiente, il task viene bloccato con notifica alla cassa; le
restanti righe proseguono normalmente.

\paragraph{Passo 4  Identificazione reparti compatibili}
Per ogni task con scorta sufficiente, il sistema individua i reparti in grado di
gestire la pietanza: il reparto primario (\texttt{settore\_stampa}) più eventuali
alternative configurate (es.\ una pietanza fritta gestibile sia da cucina che da
griglia).

\paragraph{Passo 5  Calcolo del carico per ogni reparto}
Per ciascun reparto compatibile viene calcolato un punteggio di carico:

\begin{equation}
\label{eq:carico}
\mathtt{carico\_reparto} = n_{\text{task}} \times t_{\text{prep\_residuo}} \times \alpha_{\text{stress}} \times \beta_{\text{stampante}}
\end{equation}

\noindent dove $n_{\text{task}}$ è il numero di task in coda, $t_{\text{prep\_residuo}}$
è la somma del tempo residuo di tutti i task (secondi), $\alpha_{\text{stress}} \in
[1.0,\,1.5]$ è il moltiplicatore per reparto in stress e $\beta_{\text{stampante}}
\in [1.0,\,2.0]$ è il moltiplicatore per stampante offline. Nell'implementazione,
$\alpha_{\text{stress}} = 1.3$ si attiva oltre la soglia di $5$ task in coda e
$\beta_{\text{stampante}} = 2.0$ a stampante offline.

\paragraph{Passo 6  Selezione del reparto con carico minimo}
Se il task è compatibile con più reparti, il sistema sceglie quello con
\texttt{carico\_reparto} più basso (selezione greedy). In caso di parità, prevale
il reparto con il timestamp dell'ultimo task più vecchio.

\begin{equation}
\label{eq:routing}
\mathtt{reparto\_scelto} = \arg\min_{r \,\in\, \text{reparti\_compatibili}} \mathtt{carico\_reparto}(r)
\end{equation}

Se tutti i reparti superano la soglia di saturazione, il task viene comunque
assegnato a quello meno saturo. Con un solo reparto disponibile, l'algoritmo degrada
a FIFO con priorità per massimo tempo di attesa.

\paragraph{Passo 7  Calcolo del ranking \texorpdfstring{$P$}{P}}
Una volta assegnato il task, il ranking $P$ ne determina l'ordine nella coda (valore
più alto = elaborazione anticipata):

\begin{equation}
\label{eq:ranking}
P = A \cdot t_{\text{attesa}} + B \cdot \text{priorità\_asporto} + C \cdot \text{semplicità} - D \cdot \text{carico} - E \cdot \text{rischio\_scorte}
\end{equation}

\begin{table}[H]
\centering
\caption{Termini della formula di ranking $P$}
\label{tab:ranking}
\begin{tabular}{cl}
\toprule
\textbf{Termine} & \textbf{Significato} \\
\midrule
$A \cdot t_{\text{attesa}}$             & Tempo dall'arrivo dell'ordine (urgenza crescente). \\
$B \cdot \text{priorità\_asporto}$      & Bonus per ordini asporto/urgenza con tempi vincolati. \\
$C \cdot \text{semplicità}$             & Bonus per pietanze rapide; libera il reparto nel picco. \\
$-D \cdot \text{carico}$                & Penalità per reparti già in stress. \\
$-E \cdot \text{rischio\_scorte}$       & Penalità per scorte prossime all'esaurimento. \\
\bottomrule
\end{tabular}
\end{table}

\paragraph{Passo 8  Raggruppamento per reparto e ottimizzazione batch}
I task di uno stesso ordine vengono raggruppati per reparto, producendo una comanda
separata per ciascuno. Pietanze identiche in arrivo entro la finestra di batch
(\texttt{FINESTRA\_BATCH\_MS}, pari a $30$ secondi nell'implementazione) vengono
aggregate in un unico task cumulato.

\paragraph{Passo 9  Generazione comande separate}
Per ogni reparto coinvolto viene generata una comanda indipendente. Esempio:
l'ordine \#52 (birra media, costine, patatine, tiramisù) produce tre comande:
\textsc{Bar}, \textsc{Griglia} e \textsc{Cucina}.

\paragraph{Passo 10  Ordinamento per priorità e invio}
Le comande vengono ordinate per $P$ decrescente e inviate alla stampante del
reparto. Se la stampante è offline, si usa una stampante di fallback; in assenza di
fallback, la comanda appare nella dashboard admin con notifica e richiesta di
conferma manuale.

\paragraph{Passo 11  Loop su task residui}
Il sistema verifica se restano task non processati. In caso affermativo il flusso
ritorna al Passo~7; il ciclo termina quando tutti i task dell'ordine sono stati
inviati.

\subsubsection*{Casi Limite e Comportamenti Attesi}

\begin{table}[H]
\centering
\caption{Casi limite e strategie di gestione}
\label{tab:casislimite}
\begin{tabular}{lp{4.5cm}p{5.5cm}}
\toprule
\textbf{Caso limite} & \textbf{Condizione} & \textbf{Comportamento} \\
\midrule
Stampante offline
    & La stampante del reparto non risponde.
    & Fallback su stampante secondaria o dashboard admin con notifica. \\
Reparto saturo
    & Tutti i reparti compatibili oltre la soglia.
    & Assegnazione al meno saturo; degrado a FIFO con priorità attesa. \\
Scorta insufficiente
    & Quantità disponibile $<$ richiesta.
    & Task bloccato; notifica alla cassa; altre righe non impattate. \\
Ordine misto
    & Pietanze per reparti diversi.
    & Sotto-comande indipendenti per ciascun reparto. \\
Picco improvviso
    & $N$ ordini dello stesso piatto in breve intervallo.
    & Raggruppamento batch con quantità cumulata. \\
Reparto unico
    & Un solo reparto compatibile operativo.
    & Routing forzato; FIFO con priorità per attesa massima. \\
\bottomrule
\end{tabular}
\end{table}

\subsubsection*{Analisi di Complessità}

Siano $n$ il numero di reparti, $m$ il numero di task in coda, $k$ il numero di
righe valide nell'ordine.

\begin{table}[H]
\centering
\caption{Complessità per fase}
\label{tab:complessita}
\begin{tabular}{ll}
\toprule
\textbf{Fase} & \textbf{Complessità} \\
\midrule
Lettura righe ordine                   & $O(k)$ \\
Verifica scorte (per task)             & $O(1)$ \\
Identificazione reparti compatibili    & $O(n)$ \\
Calcolo \texttt{carico\_reparto}       & $O(n \cdot m)$ \\
Selezione $\min(\text{carico})$        & $O(n)$ \\
Calcolo ranking $P$                    & $O(1)$ \\
Ordinamento comande                    & $O(k \log k)$ \\
Invio comanda (per reparto)            & $O(1)$ \\
\midrule
\textbf{Totale per ordine}             & $\boldsymbol{O(n \cdot m + k \log k)}$ \\
\bottomrule
\end{tabular}
\end{table}

In pratica $n \leq 10$ e $k \leq 20$: l'algoritmo è estremamente efficiente. Il
termine dominante è $O(n \cdot m)$, dove $m$ dipende dalla lunghezza delle code a
runtime.

\subsubsection*{Pseudocodice}

Le Figure~\ref{alg:routeorder} e~\ref{alg:subfunzioni} mostrano lo pseudocodice
completo.

\begin{figure}[H]
\begin{algorithmic}[1]
\Procedure{RouteOrder}{$\text{order}$}
    \State $\text{tasks} \gets [\,]$
    \For{$\text{riga} \in \text{order.righe}$} \Comment{$O(k)$}
        \State $\text{task} \gets \Call{BuildTask}{\text{riga},\;\text{order.timestamp},\;\text{order.priorità}}$
        \State $\text{tasks}.\text{append}(\text{task})$
    \EndFor
    \For{$\text{task} \in \text{tasks}$} \Comment{$O(k)$}
        \If{$\neg\Call{ScorteDisponibili}{\text{task}}$}
            \State $\Call{Blocca}{\text{task}}$;\; $\Call{NotificaCassa}{\text{task}}$
            \State $\text{tasks}.\text{remove}(\text{task})$
        \EndIf
    \EndFor
    \State $\text{comande} \gets \{\}$
    \For{$\text{task} \in \text{tasks}$} \Comment{$O(k \cdot n \cdot m)$}
        \State $\text{reparti} \gets \Call{GetRepartiCompatibili}{\text{task}}$ \Comment{$O(n)$}
        \State $r \gets \Call{SelectMinCarico}{\text{reparti}}$ \Comment{$O(n \cdot m)$}
        \State $\text{task}.P \gets \Call{CalcolaRanking}{\text{task},\;r}$ \Comment{$O(1)$}
        \State $\text{comande}[r].\text{append}(\text{task})$
    \EndFor
    \State $\text{comande} \gets \Call{RaggruppaEBatch}{\text{comande}}$ \Comment{$O(k \log k)$}
    \For{$r \in \Call{Keys}{\text{comande}}$}
        \State $\text{task\_list} \gets \Call{SortByP}{\text{comande}[r]}$ \Comment{$O(k \log k)$}
        \State $\text{comanda} \gets \Call{GeneraComanda}{r,\;\text{task\_list}}$
        \State $\Call{InviaComanda}{\text{comanda},\;r}$
    \EndFor
\EndProcedure
\end{algorithmic}
\caption{Pseudocodice funzione principale \textsc{RouteOrder}}
\label{alg:routeorder}
\end{figure}

\begin{figure}[H]
\begin{algorithmic}[1]
\Procedure{SelectMinCarico}{$\text{reparti}$} \Comment{$O(n \cdot m)$}
    \State $\text{best} \gets \textbf{null}$;\quad $\text{best\_score} \gets +\infty$
    \For{$r \in \text{reparti}$}
        \State $\text{score} \gets \displaystyle\sum_{\text{task} \in r.\text{coda}} \text{task}.t_{\text{prep\_residuo}}$
        \State $\text{score} \gets \text{score} \times \alpha_{\text{stress}}(r) \times \beta_{\text{stampante}}(r)$
        \If{$\text{score} < \text{best\_score}$}
            \State $\text{best} \gets r$;\; $\text{best\_score} \gets \text{score}$
        \EndIf
    \EndFor
    \State \Return $\text{best}$
\EndProcedure
\Statex
\Procedure{CalcolaRanking}{$\text{task},\;r$} \Comment{$O(1)$}
    \State $t_{\text{attesa}} \gets \Call{Now}{} - \text{task.timestamp}$
    \State $\pi_a \gets [\text{task.priorità} = \textsc{Asporto}]\;?\;1.0 : 0.0$
    \State $s \gets 1.0\;/\;(\text{task}.t_{\text{prep}} + 1)$
    \State \Return $A \cdot t_{\text{attesa}} + B \cdot \pi_a + C \cdot s - D \cdot \text{carico}(r) - E \cdot \text{rischio\_scorte}(\text{task})$
\EndProcedure
\end{algorithmic}
\caption{Pseudocodice funzioni ausiliarie}
\label{alg:subfunzioni}
\end{figure}

\clearpage
\begin{figure}[H]
    \centering
    \includegraphics[height=0.90\textheight, keepaspectratio]{smistatore.png}
    \caption{Diagramma di flusso  Smistatore Intelligente Multi-Reparto}
    \label{fig:flusso_smistatore}
\end{figure}

\clearpage
% ============================================================================
\subsection{Test}

\subsubsection{Analisi statica del codice}

Il backend è stato sottoposto ad analisi statica con gli stessi strumenti usati nel
resto del progetto: la \textbf{complessità ciclomatica} delle funzioni è stata
misurata con la regola \texttt{complexity} di \textbf{ESLint}, mentre
l'\textbf{accoppiamento} tra moduli e le dipendenze circolari sono stati analizzati
con \textbf{madge}. Per \emph{isolare} il contributo dello Smistatore alla
complessità del sistema, l'analisi confronta il backend \emph{base} (privo di
entrambi gli algoritmi) con lo stesso backend a cui è stato aggiunto il \emph{solo}
Smistatore, tenendo temporaneamente fuori il Predittore: la differenza tra i due
stati è dunque attribuibile interamente allo Smistatore.

\begin{table}[H]
\centering
\caption{Metriche statiche del backend: base contro stato con il solo Smistatore}
\label{tab:analisi-smistatore}
\begin{tabular}{lccc}
\toprule
\textbf{Metrica} (backend, \texttt{src/}) & \textbf{Base} & \textbf{Con Smistatore} & \textbf{$\Delta$} \\
\midrule
Complessità ciclomatica massima        & 51 & 51 & $=$ \\
Funzioni con complessità ciclomatica ${}>{}$ 15 & 4 & 4 & $=$ \\
Funzioni totali analizzate             & 67 & 85 & $+18$ \\
Moduli accoppiati a \texttt{db.js} ($C_a$) & 7 & 8 & $+1$ \\
Dipendenze circolari                   & 0 & 0 & $=$ \\
Righe di codice (SLOC, escluse vuote e commenti) & 1\,506 & 1\,732 & $+226$ \\
\bottomrule
\end{tabular}
\end{table}

% "parsaFlusso" e' il nome REALE della funzione in services/escpos-emulator.js.
Lo Smistatore introduce 18 nuove funzioni e 226 righe di codice, senza generare
dipendenze circolari. Tutte le sue funzioni restano a bassa complessità: la più
impegnativa è \texttt{rigaToTasks} (15), che converte le righe d'ordine in task,
mentre la funzione principale di routing \texttt{routeOrder} si ferma a 6. La
complessità ciclomatica massima del sistema resta invariata (51, relativa al parser
ESC/POS \texttt{parsaFlusso} nell'emulatore, estraneo allo Smistatore): l'algoritmo
non peggiora cioè il punto più critico del codice, e nessuna delle funzioni
introdotte supera la soglia di 15. La complessità $51$ di \texttt{parsaFlusso} è un
debito tecnico noto, accettato in questa fase perché il parser è stabile e isolato,
e affrontato con un refactoring dedicato nell'Iterazione~3. I tre grafici seguenti
visualizzano l'impatto dello Smistatore.

\definecolor{statBase}{HTML}{D45454}
\definecolor{statAlg}{HTML}{5BA85E}
\pgfplotsset{
    statbar/.style={
        xbar, width=\linewidth, bar width=7pt, enlarge y limits=0.22, xmin=0,
        axis x line*=bottom, axis y line*=left, ytick=data,
        nodes near coords, nodes near coords align={horizontal},
        every node near coord/.append style={font=\scriptsize},
        legend style={at={(0.97,0.05)}, anchor=south east, font=\small, draw=none},
        legend cell align={left}, tick label style={font=\small}, xlabel style={font=\small},
    }
}

\begin{figure}[H]
    \centering
    \resizebox{\textwidth}{!}{%
    \begin{tikzpicture}
        \begin{axis}[statbar, height=6cm, xmax=58,
            xlabel={Complessità ciclomatica (CC)},
            symbolic y coords={p1,p2,p3,p4,p5},
            yticklabels={routeOrder, batchPerReparto, rigaToTasks,
                         avviaServer, parsaFlusso},
        ]
            \addplot[fill=statBase, draw=statBase] coordinates {(0,p1)(0,p2)(0,p3)(31,p4)(51,p5)};
            \addplot[fill=statAlg, draw=statAlg]   coordinates {(6,p1)(7,p2)(15,p3)(31,p4)(51,p5)};
            \legend{Base, Con Smistatore}
        \end{axis}
    \end{tikzpicture}%
    }
    \caption{Complessità ciclomatica delle funzioni principali del backend. Le tre
    funzioni in basso (\texttt{routeOrder}, \texttt{batchPerReparto},
    \texttt{rigaToTasks}) sono introdotte dallo Smistatore e restano tutte sotto la
    soglia di 15; le altre, legate all'infrastruttura di base, sono invariate.}
    \label{fig:cc-smistatore}
\end{figure}

\begin{figure}[H]
    \centering
    \resizebox{\textwidth}{!}{%
    \begin{tikzpicture}
        \begin{axis}[statbar, height=4.2cm, xmax=14,
            xlabel={Conteggio},
            symbolic y coords={s1,s2,s3},
            yticklabels={Funzioni con CC\,$>$\,15, Moduli totali,
                         Moduli accoppiati al DB},
        ]
            \addplot[fill=statBase, draw=statBase] coordinates {(4,s1)(10,s2)(7,s3)};
            \addplot[fill=statAlg, draw=statAlg]   coordinates {(4,s1)(12,s2)(8,s3)};
            \legend{Base, Con Smistatore}
        \end{axis}
    \end{tikzpicture}%
    }
    \caption{Metriche strutturali prima e dopo l'aggiunta del solo Smistatore: il
    modulo aggiunge due moduli al grafo (route e service), di cui uno (la route)
    accoppiato al database.}
    \label{fig:strutt-smistatore}
\end{figure}

\begin{figure}[H]
    \centering
    \resizebox{\textwidth}{!}{%
    \begin{tikzpicture}
        \begin{axis}[statbar, height=3.4cm, xmax=1850,
            xlabel={Righe di codice (SLOC)},
            symbolic y coords={d1},
            yticklabels={Backend (\texttt{src/})},
        ]
            \addplot[fill=statBase, draw=statBase] coordinates {(1506,d1)};
            \addplot[fill=statAlg, draw=statAlg]   coordinates {(1732,d1)};
            \legend{Base, Con Smistatore}
        \end{axis}
    \end{tikzpicture}%
    }
    \caption{Dimensione del backend in righe di codice (SLOC) prima e dopo
    l'aggiunta del solo Smistatore.}
    \label{fig:sloc-smistatore}
\end{figure}

\subsubsection{Test automatizzati}

Per validare la correttezza dell'algoritmo è stata sviluppata una suite di test
automatizzati con \textbf{Vitest}, eseguibile con \texttt{npm test}. La suite
comprende \textbf{test unitari}, che esercitano in isolamento le funzioni pure di
calcolo del modulo, e \textbf{test di integrazione HTTP}, che verificano il
comportamento end-to-end sulle API REST con il backend in esecuzione (saltati
automaticamente quando il server non è raggiungibile).{fig:test-smistatore}.

\begin{figure}[H]
    \centering
    \includegraphics[width=0.92\textwidth]{IT2_test_smistatore.png}
    \caption{Esecuzione della suite unitaria dello Smistatore con Vitest:
             26/26 test unitari superati}
    \label{fig:test-smistatore}
\end{figure}

I test unitari sono organizzati per componente dell'algoritmo:

\begin{enumerate}
    \item \textbf{Formula del carico.} Verificano il calcolo $n \times t_{prep}
    \times \alpha \times \beta$, con l'attivazione del fattore di stress
    $\alpha=1.3$ oltre la soglia di 5 task e del fattore $\beta=2.0$ a stampante
    offline.

    \item \textbf{Selezione del reparto} (\texttt{selectMinCarico}). Validano la
    scelta greedy del reparto a carico minore, il tie-break sul task più vecchio a
    parità di carico e il fallback quando tutte le stampanti sono offline.

    \item \textbf{Routing e batching} (\texttt{routeOrder}). Coprono la generazione
    di un task per riga, il blocco delle righe a scorta insufficiente (mentre le
    altre proseguono), l'esplosione delle modalità \emph{singola\_singola} /
    \emph{singola\_multipla} e il raggruppamento di più voci dello stesso reparto in
    un'unica comanda.

    \item \textbf{Ranking $P$.} Verificano il bonus di priorità per gli ordini da
    asporto ($B=50$) e la penalizzazione per il rischio di esaurimento scorte
    ($E=30$).

    \item \textbf{Bilanciamento e fallback.} Validano la distribuzione del carico su
    reparti con fallback bidirezionale (6 ordini ripartiti 3-3) e il re-routing
    automatico verso il reparto alternativo quando la stampante primaria è offline.
\end{enumerate}


\subsubsection{API esposte}

Gli endpoint seguenti supportano lo Smistatore e la catena di stampa; valgono le
convenzioni REST definite nell'Iterazione~1. Le operazioni di scrittura
(\POST, \PUT) sono verificate tramite Postman, di cui si riporta richiesta e
risposta.

\paragraph{Smistatore}\mbox{}\\

Espone lo stato interno dello \textbf{Smistatore Intelligente Multi-Reparto}
(UC7) e permette di aggiornare a runtime le configurazioni di fallback e lo stato
delle stampanti, senza riavviare il server.

\apicall{\GET~\texttt{/api/smistatore/stato}}

Snapshot delle code di preparazione per ogni reparto attivo: numero di task in coda,
carico stimato e stato della stampante associata.

\begin{lstlisting}[style=json]
{
  "cucina":  { "task_in_coda": 3, "carico": 420, "stato_stampante": "online"  },
  "bar":     { "task_in_coda": 1, "carico":  60, "stato_stampante": "online"  },
  "griglia": { "task_in_coda": 0, "carico":   0, "stato_stampante": "offline" }
}
\end{lstlisting}

\apicall{\POST~\texttt{/api/smistatore/completa}}

Rimuove un task dalla coda di un reparto una volta che la preparazione è completata.
Campi obbligatori: \texttt{reparto} e \texttt{voce\_id}. Campo opzionale:
\texttt{ordine\_id} per identificare univocamente il task in caso di omonimi.

\postman{i2-smist-completa.png}{Postman  \texttt{POST /api/smistatore/completa}}

\apicall{\PUT~\texttt{/api/smistatore/stampante/:reparto}}

Imposta lo stato (\texttt{online} / \texttt{offline}) della stampante di un reparto,
aggiornando sia la memoria interna dello smistatore che il database. Se la stampante
passa a \texttt{offline}, i task successivi vengono automaticamente instradati al
fallback configurato.

\postman{i2-smist-stampante.png}{Postman  \texttt{PUT /api/smistatore/stampante/:reparto}}

\apicall{\PUT~\texttt{/api/smistatore/fallback/:reparto}}

Configura i reparti alternativi a cui instradare i task quando la stampante del
reparto primario è offline. Il campo \texttt{fallback} è un array di nomi reparto.

\postman{i2-smist-fallback.png}{Postman  \texttt{PUT /api/smistatore/fallback/:reparto}}

\paragraph{Stampe}\mbox{}\\

Gestisce lo storico delle stampe ESC/POS inviate, la ristampa di comande e il
controllo dell'emulatore TCP usato in modalità demo/sviluppo.

\apicall{\GET~\texttt{/api/stampe}}

Storico delle stampe eseguite. Con il parametro \texttt{?ordine\_id=N} restituisce
solo le stampe relative a quell'ordine; senza parametro restituisce le ultime 100 in
ordine cronologico decrescente.

\begin{lstlisting}[style=json]
[
  { "id": 1, "ordine_id": 42, "reparto": "cucina", "esito": "ok",
    "durata_ms": 85, "timestamp": "2026-05-26T20:14:00.000Z",
    "stampante_nome": "Cucina A", "indirizzo_ip": "127.0.0.1", "porta": 9100 }
]
\end{lstlisting}

\apicall{\GET~\texttt{/api/stampe/recenti}}

Ultime $N$ stampe eseguite (default 50, massimo 500) con parametro opzionale
\texttt{?limit=N}. Versione alleggerita senza payload per uso nelle dashboard.

\begin{lstlisting}[style=json]
[
  { "id": 5, "ordine_id": 44, "reparto": "bar", "esito": "ok",
    "durata_ms": 42, "timestamp": "2026-05-26T20:15:10.000Z",
    "stampante_nome": "Stampante Bar" }
]
\end{lstlisting}

\apicall{\POST~\texttt{/api/stampe/:id/ristampa}}

Reinvia alla stampante la comanda corrispondente a una stampa già registrata,
riutilizzando il payload JSON salvato. Risposta 400 se il payload non è disponibile;
404 se la stampa non esiste.

\begin{lstlisting}[style=json]
{ "ristampa": { "esito": "ok", "durata_ms": 91 } }
\end{lstlisting}

\apicall{\POST~\texttt{/api/stampe/test/:id}}

Invia una pagina di test alla stampante identificata da \texttt{id}. Utile per
verificare la connettività prima dell'evento.

\begin{lstlisting}[style=json]
{ "esito": "ok", "stampante_id": 2, "durata_ms": 120 }
\end{lstlisting}

\apicall{\POST~\texttt{/api/stampe/emulatore/:stampante\_id/spegni}}

Spegne la stampante emulata indicata, aggiorna lo stato nel database a
\texttt{offline} e notifica lo smistatore tramite WebSocket
(\texttt{stampante\_offline}). Usato per simulare guasti in modalità demo.

\begin{lstlisting}[style=json]
{ "ok": true, "stampante_id": 1, "stato": "offline" }
\end{lstlisting}

\apicall{\POST~\texttt{/api/stampe/emulatore/:stampante\_id/accendi}}

Riavvia la stampante emulata, aggiorna lo stato a \texttt{online} e notifica lo
smistatore tramite WebSocket (\texttt{stampante\_online}).

\begin{lstlisting}[style=json]
{ "ok": true, "stampante_id": 1, "stato": "online" }
\end{lstlisting}

\apicall{\GET~\texttt{/api/stampe/emulatore/stato}}

Snapshot dello stato di tutti gli emulatori TCP attivi (porta, stato, ultima
attività).

\begin{lstlisting}[style=json]
{
  "1": { "porta": 9100, "stato": "online",  "ultima_attivita": "2026-05-26T20:14:00Z" },
  "2": { "porta": 9101, "stato": "online",  "ultima_attivita": "2026-05-26T20:12:30Z" },
  "3": { "porta": 9102, "stato": "offline", "ultima_attivita": null }
}
\end{lstlisting}

\newpage